\chapter[Autre famille de modèles, seuils et informations absentes]{Autre famille de modèles,\\seuils et informations absentes}
\label{chap:comparaison}

Le chapitre~\ref{chap:test-ml} se termine sur trois réserves explicites~: le modèle additif
ne capture pas toutes les relations possibles, le rappel au seuil illustratif reste faible
et l'absence d'information porte un signal particulier. Le notebook 08 reprend ces trois
points. Il compare une seconde famille de modèles, rend l'arbitrage des seuils explicite
et documente le groupe des usagers sans âge.

\section{Ce que ce chapitre peut et ne peut pas établir}
Les scores de 2024 de la régression logistique ont été consultés au chapitre précédent.
Cette année n'est donc plus un jeu resté inconnu pour la comparaison menée ici. Deux
garde-fous sont appliqués et doivent être rappelés avec les résultats~:

\begin{itemize}\setlength{\itemsep}{1pt}
\item les réglages du modèle d'arbres sont choisis \textbf{sur 2023 seule}, avec une
grille écrite sur disque avant le premier entraînement~;
\item aucun score de 2024 n'intervient dans ce choix, et aucun réentraînement n'a lieu
après lecture du test.
\end{itemize}

La stabilité observée entre 2023 et 2024 rend le résultat crédible. Elle ne remplace pas
une année réellement vierge, qui serait nécessaire pour trancher définitivement entre les
deux familles. Le protocole des chapitres~\ref{chap:protocole-ml} à~\ref{chap:test-ml}
n'est pas modifié rétrospectivement~: ce chapitre ajoute une comparaison, il ne révise
aucun critère antérieur.

\section{Reconstruire avant de comparer}
Comme le notebook 07, le notebook 08 vérifie l'empreinte SHA-256 de la table préparée,
puis réentraîne la régression logistique retenue et compare ses scores de 2023 à ceux
enregistrés par le notebook 06. L'écart maximal observé reste inférieur à la tolérance
de $10^{-4}$. Sans ce contrôle, un écart entre deux modèles pourrait provenir d'un
changement d'entrée plutôt que de la méthode comparée.

\section{Un protocole d'arbres fixé avant les scores}
Le modèle comparé est un ensemble d'arbres de décision construits séquentiellement par
gradient boosting~\cite{sklearn-hgb}. Chaque arbre corrige les erreurs des précédents,
et la somme de leurs contributions est transformée en probabilité. Contrairement à la
régression logistique, cette construction peut représenter une combinaison de
caractéristiques dont l'effet diffère de la somme de leurs effets séparés.

Les onze variables, les tranches d'âge et le découpage temporel sont identiques à ceux
du chapitre~\ref{chap:protocole-ml}. Deux éléments seulement changent, et chacun est
justifié~:

\textbf{La représentation des modalités.} L'implémentation traite nativement les
variables catégorielles, sans encodage indicateur. Les modalités sont apprises sur
2019-2022 et déclarées comme telles~; une modalité jamais vue devient une valeur
manquante, gérée comme telle et comptée. Aucune modalité nouvelle n'apparaît en 2023
ni en 2024.

\textbf{L'arrêt anticipé est désactivé.} Son fonctionnement habituel consiste à mettre
de côté une fraction aléatoire de l'apprentissage pour décider quand cesser d'ajouter
des arbres. Ce découpage séparerait les usagers d'un même accident, alors que tout le
projet les maintient groupés. Le nombre d'itérations est donc placé dans la grille et
choisi sur 2023, comme les autres réglages.

La grille comporte trois réglages à deux valeurs chacun~: le pas d'apprentissage
(0,05 ou 0,1), le nombre maximal de feuilles par arbre (15 ou 31) et le nombre
d'itérations (200 ou 400). Les autres paramètres sont fixés~: au moins 50 usagers par
feuille, aucune pénalisation L2 supplémentaire et une graine identique à celle du
notebook 06. La règle de sélection reprend celle du chapitre~\ref{chap:validation-ml}~:
log-loss minimale sur 2023, égalité départagée par le modèle le plus simple, et modèle
conservé seulement s'il fait strictement mieux.

\section{Sélection sur 2023}
\begin{table}[htbp]
\centering\small
\renewcommand{\arraystretch}{1.18}
\caption{Les huit configurations d'arbres évaluées sur 2023 ; ligne retenue en gras.}\label{tab:ml-selection-arbres}
\begin{tabularx}{\textwidth}{@{}rrrrrrr@{}}
\toprule
Pas & Feuilles & Itérations & Log-loss & Brier & AP & ROC-AUC \\
\midrule
0,05 & 15 & 200 & 0,386638 & 0,122034 & 0,456584 & 0,787240 \\
0,05 & 15 & 400 & 0,385326 & 0,121621 & 0,460178 & 0,788943 \\
0,05 & 31 & 200 & 0,384799 & 0,121441 & 0,461567 & 0,789662 \\
\textbf{0,05} & \textbf{31} & \textbf{400} & \textbf{0,384591} & \textbf{0,121401} & \textbf{0,462051} & \textbf{0,789973} \\
0,10 & 15 & 200 & 0,385420 & 0,121657 & 0,459306 & 0,788906 \\
0,10 & 15 & 400 & 0,385012 & 0,121545 & 0,460528 & 0,789571 \\
0,10 & 31 & 200 & 0,384906 & 0,121518 & 0,460611 & 0,789676 \\
0,10 & 31 & 400 & 0,385496 & 0,121744 & 0,458971 & 0,789144 \\
\bottomrule
\end{tabularx}
\sourcecalcul{notebook 08 ; grille écrite avant le premier entraînement, sélection par log-loss de 2023}
\end{table}


La configuration retenue utilise un pas de 0,05, 31 feuilles et 400 itérations.
Les huit candidats tiennent toutefois dans un intervalle de log-loss d'environ 0,002,
entre 0,3846 et 0,3866.
Comme pour les quatre valeurs de $C$ au chapitre~\ref{chap:validation-ml}, le réglage
n'est pas déterminant~: c'est le choix de la famille de modèles qui produit l'écart,
non son accordage. Présenter cette sélection comme une optimisation décisive serait
trompeur.

\section{Résultats sur 2024}
\begin{table}[htbp]
\centering\small
\renewcommand{\arraystretch}{1.18}
\caption{Référence, régression logistique et arbres boostés sur 2023 et 2024.}\label{tab:ml-comparaison}
\begin{tabularx}{\textwidth}{@{}lLrrrr@{}}
\toprule
Année & Modèle & Log-loss & Brier & AP & ROC-AUC \\
\midrule
2023 & Référence constante & 0,4709 & 0,1474 & 0,1797 & 0,5000 \\
2023 & Régression logistique & 0,4008 & 0,1271 & 0,4085 & 0,7699 \\
\textbf{2023} & \textbf{Arbres boostés} & \textbf{0,3846} & \textbf{0,1214} & \textbf{0,4621} & \textbf{0,7900} \\
2024 & Référence constante & 0,4703 & 0,1471 & 0,1793 & 0,5000 \\
2024 & Régression logistique & 0,4001 & 0,1268 & 0,4089 & 0,7699 \\
\textbf{2024} & \textbf{Arbres boostés} & \textbf{0,3845} & \textbf{0,1213} & \textbf{0,4629} & \textbf{0,7890} \\
\bottomrule
\end{tabularx}
\sourcecalcul{notebook 08 ; même apprentissage 2019-2022 et même représentation pour les deux modèles}
\end{table}


Le modèle d'arbres améliore les quatre indicateurs, sur 2023 comme sur 2024. Sur le test,
la log-loss passe de 0,4001 à 0,3845 et le score de Brier de 0,1268 à 0,1213. Rapportée
à la référence constante, la réduction de log-loss passe ainsi de 14,93\,\% pour la
logistique à 18,24\,\% pour les arbres.

Deux lectures méritent d'être séparées. D'une part, le gain en log-loss reste modeste~:
environ 3,9\,\% par rapport à la logistique. D'autre part, l'\textit{average precision}
progresse beaucoup plus nettement, de 0,4089 à 0,4629, soit 13,2\,\%. Ce second
indicateur se concentre sur le haut du classement. L'écart entre les deux progressions
indique que l'amélioration ne se répartit pas uniformément~: elle porte surtout sur
l'identification des usagers les plus exposés.

La stabilité temporelle se vérifie pour les deux modèles. La ROC-AUC de la logistique
vaut 0,7699 en 2023 et 0,7699 en 2024~; celle des arbres 0,7900 et 0,7890. Ce qui est
appris sur 2019-2022 conserve le même niveau de performance deux ans plus tard. Une
seule année supplémentaire ne suffit cependant pas à établir une stabilité générale.

\section{Pourquoi les arbres font mieux}
L'écart de score ne dit pas d'où il vient. Deux diagnostics le précisent.

\subsection{La logistique plafonne dans le haut du risque}
\figurerapportML[0.88]{08_01_calibration_comparaison_test.png}{Calibration des deux modèles sur 2024, dix intervalles fixes. Les intervalles de moins de 50 usagers sont écartés.}{fig:calibration-comparaison}{08}

En dessous de 0,45, les deux courbes suivent la diagonale. Au-delà, la régression
logistique s'en écarte~: le chapitre~\ref{chap:test-ml} relevait déjà qu'elle annonce
74,28\,\% dans l'intervalle $]0{,}7;0{,}8]$ pour 59,90\,\% observés. Elle n'attribue
par ailleurs presque jamais de probabilité supérieure à 0,85. Les arbres restent proches
de la diagonale sur toute l'étendue et atteignent des valeurs plus élevées.

Autrement dit, le modèle additif estime correctement le risque moyen mais \textbf{ne sait
pas désigner les cas les plus graves}~: il surestime les probabilités élevées tout en
étant incapable d'en produire de très élevées. C'est précisément ce que mesure
l'\textit{average precision}, et cela explique pourquoi elle progresse davantage que
la log-loss.

\subsection{Les combinaisons de caractéristiques}
\begin{table}[htbp]
\centering\small
\renewcommand{\arraystretch}{1.18}
\caption{Profils les plus exposés en 2024 : proportions observées et probabilités moyennes.}\label{tab:ml-profils}
\begin{tabularx}{\textwidth}{@{}LLrrrr@{}}
\toprule
Catégorie & Tranche d'âge & Usagers & Observé (\%) & Logistique (\%) & Arbres (\%) \\
\midrule
Piéton & 75 ans ou plus & 1\,500 & 50,67 & 49,00 & 52,12 \\
Piéton & 65-74 ans & 930 & 41,18 & 40,11 & 41,74 \\
Piéton & 55-64 ans & 931 & 34,26 & 32,98 & 34,44 \\
Passager & 75 ans ou plus & 919 & 34,17 & 30,72 & 33,58 \\
Piéton & 45-54 ans & 970 & 32,27 & 28,85 & 29,83 \\
Conducteur & 0-17 ans & 2\,993 & 31,94 & 25,53 & 31,34 \\
Piéton & 35-44 ans & 919 & 28,94 & 26,95 & 27,97 \\
Piéton & 18-24 ans & 844 & 26,30 & 29,44 & 26,22 \\
\bottomrule
\end{tabularx}
\sourcecalcul{notebook 08 ; groupes d'au moins 300 usagers, triés par proportion observée}
\end{table}


Sur les profils croisant catégorie d'usager et tranche d'âge, restreints aux groupes
d'au moins 300 usagers, l'écart absolu moyen entre probabilité moyenne estimée et
proportion observée vaut 1,52 point pour la régression logistique et 0,55 point pour
les arbres.

Le cas des \textbf{conducteurs de moins de 18 ans} est le plus parlant~: pour la plupart
des conducteurs de deux-roues motorisés, compte tenu de l'âge légal de conduite. La
proportion observée y atteint 31,94\,\% en 2024, alors que la régression logistique
annonce 25,53\,\%. Sur l'échelle des log-cotes, elle additionne un effet de tranche
d'âge modéré et un effet de catégorie neutre~; elle ne dispose d'aucun terme pour
exprimer que la combinaison des deux dépasse la somme de ses parties. Les arbres
retrouvent 31,34\,\%.

Ce diagnostic décrit un défaut de représentation, pas un effet causal de l'âge ou du
statut de conducteur. Il indique où le modèle additif se trompe, sans établir pourquoi
ces usagers sont plus exposés.

\section{Rendre l'arbitrage des seuils explicite}
Le seuil de 0,5 des chapitres précédents est illustratif. Il conduit à un rappel de
18,22\,\% pour la logistique et de 21,16\,\% pour les arbres. Ces valeurs ne disent rien
du coût d'un rappel plus élevé, et les lire comme une limite du modèle serait une erreur.

Fixer le seuil directement sur 2024 reviendrait à régler une quantité sur l'année de
test. Chaque seuil est donc déterminé sur 2023, pour y atteindre un rappel visé, puis
appliqué sans retouche à 2024.
\begin{table}[htbp]
\centering\small
\renewcommand{\arraystretch}{1.18}
\caption{Seuils fixés sur 2023 pour un rappel visé, puis appliqués sans retouche à 2024.}\label{tab:ml-seuils}
\begin{tabularx}{\textwidth}{@{}Lrrrrrrr@{}}
\toprule
Modèle & Visé (\%) & Seuil & Rappel (\%) & Précision (\%) & Fausses alertes & Cas manqués & Alertés (\%) \\
\midrule
Logistique & 20 & 0,477 & 20,14 & 52,98 & 3\,757 & 16\,790 & 6,81 \\
Logistique & 30 & 0,377 & 30,12 & 47,48 & 7\,004 & 14\,691 & 11,37 \\
Logistique & 40 & 0,307 & 40,25 & 43,21 & 11\,123 & 12\,561 & 16,70 \\
Logistique & 50 & 0,258 & 50,16 & 39,90 & 15\,883 & 10\,479 & 22,53 \\
Logistique & 60 & 0,214 & 60,47 & 36,72 & 21\,910 & 8\,310 & 29,52 \\
Logistique & 70 & 0,175 & 70,06 & 32,93 & 30\,002 & 6\,295 & 38,14 \\
Arbres & 20 & 0,514 & 19,85 & 62,78 & 2\,475 & 16\,850 & 5,67 \\
Arbres & 30 & 0,415 & 29,92 & 55,00 & 5\,147 & 14\,733 & 9,75 \\
Arbres & 40 & 0,348 & 39,94 & 48,23 & 9\,015 & 12\,626 & 14,85 \\
Arbres & 50 & 0,298 & 49,71 & 43,12 & 13\,786 & 10\,574 & 20,67 \\
Arbres & 60 & 0,244 & 60,06 & 38,93 & 19\,803 & 8\,398 & 27,65 \\
Arbres & 70 & 0,198 & 69,64 & 35,03 & 27\,150 & 6\,383 & 35,63 \\
\bottomrule
\end{tabularx}
\sourcecalcul{notebook 08 ; « Logistique » et « Arbres » désignent les deux modèles du tableau~\ref{tab:ml-comparaison} ; toutes les colonnes de résultat portent sur 2024}
\end{table}

\figurerapportML[0.86]{08_02_precision_rappel_comparaison_test.png}{Courbes précision-rappel sur 2024. Les points marquent les seuils fixés sur 2023.}{fig:pr-comparaison}{08}

Deux enseignements se dégagent.

\textbf{Les seuils se transportent.} Un seuil réglé sur 2023 pour un rappel de 50\,\%
donne un rappel de 49,7\,\% à 50,2\,\% sur 2024 selon le modèle. La stabilité temporelle
ne concerne donc pas seulement les scores agrégés~: elle vaut aussi pour une règle de
décision concrète. Ce constat porte sur deux années consécutives et sur cette population.

\textbf{L'arbitrage devient lisible.} Avec le modèle d'arbres, retrouver la moitié des
cas graves de 2024 suppose de signaler 20,7\,\% des usagers, avec 43,1\,\% d'alertes
justes, soit environ 1,3 fausse alerte par cas détecté. Les arbres dominent la
régression logistique à chaque niveau de rappel.

Ces chiffres décrivent un tableau de classification. Ils ne désignent pas un dispositif
de secours et ne définissent pas un seuil «~correct~». Le choix d'un seuil demanderait
un objectif, des coûts d'erreur et une population d'usage, qui ne sont pas établis ici.

\section{Les usagers sans âge : une pratique d'enregistrement datée}
Le chapitre~\ref{chap:exploration} signalait que les usagers sans âge à gravité connue
de 2019-2022 ne comptent que 10 décès ou hospitalisations sur 5\,550. Ce constat se
précise maintenant sur les six années.
\begin{table}[htbp]
\centering\small
\renewcommand{\arraystretch}{1.18}
\caption{Usagers sans âge et gravité associée, par année, sur les gravités connues.}\label{tab:ages-manquants}
\begin{tabularx}{\textwidth}{@{}lrrrrrrr@{}}
\toprule
Année & Usagers & Sans âge & Part (\%) & Dont nettoyage & Positifs & Taux sans âge (\%) & Taux avec âge (\%) \\
\midrule
2019 & 126\,352 & 162 & 0,13 & 162 & 9 & 5,56 & 18,10 \\
2020 & 99\,169 & 10 & 0,01 & 10 & 1 & 10,00 & 18,29 \\
2021 & 122\,135 & 2\,860 & 2,34 & 0 & 0 & 0,00 & 17,44 \\
2022 & 119\,423 & 2\,518 & 2,11 & 0 & 0 & 0,00 & 18,21 \\
2023 & 118\,101 & 2\,350 & 1,99 & 0 & 0 & 0,00 & 18,33 \\
2024 & 117\,278 & 2\,419 & 2,06 & 0 & 1 & 0,04 & 18,30 \\
\bottomrule
\end{tabularx}
\sourcecalcul{notebook 08 ; la colonne « dont nettoyage » compte les âges rendus manquants par le notebook 03}
\end{table}

\figurerapportML[0.98]{08_03_ages_manquants.png}{Fréquence de l'absence d'âge et gravité associée, par année.}{fig:ages-manquants}{08}

\subsection{Une rupture en 2021}
En 2019 et 2020, les seuls âges manquants sont les 162 et 10 valeurs que le notebook 03
a lui-même rendues manquantes parce qu'elles sortaient de l'intervalle 0-110 ans.
À partir de 2021, environ 2\,\% des usagers arrivent sans âge dans la source, soit
entre 2\,350 et 2\,860 personnes par an, et \textbf{aucun ne provient de notre
nettoyage}. La rupture tient donc à la collecte, pas à la préparation.

Cette date correspond à un changement documenté. L'audit du chapitre~\ref{chap:audit}
relevait, d'après la documentation ONISR, l'ajout d'usagers en fuite à partir de 2021,
avec des informations parfois manquantes~\cite{doc-onisr}. Le profil de ce groupe est
cohérent avec cette explication~: sur les 10\,319 usagers sans âge à gravité connue de
2019 à 2024, \textbf{99,5\,\% sont des conducteurs}, \textbf{93,8\,\% n'ont pas non plus
de sexe renseigné} et \textbf{99,4\,\% sont enregistrés comme indemnes}.

La concordance entre la date, le profil et la documentation rend cette interprétation
plausible. Elle ne la démontre pas~: les CSV ne comportent aucun champ désignant la
fuite, et rien ne permet d'attribuer individuellement chaque absence à cette cause.
Ces usagers ne sont donc pas supprimés.

\subsection{Ce que ce groupe apporte aux scores}
Le taux de décès ou d'hospitalisation y est de 0\,\% à 0,04\,\% selon l'année, contre
environ 18\,\% ailleurs. Les deux modèles l'apprennent et attribuent à ce groupe une
probabilité moyenne de l'ordre de 0,002 à 0,003 en 2024, pour 0,04\,\% observés.
\begin{table}[htbp]
\centering\small
\renewcommand{\arraystretch}{1.18}
\caption{Scores de 2024 avec et sans les usagers dont l'âge est absent.}\label{tab:ages-impact}
\begin{tabularx}{\textwidth}{@{}LLrrrrr@{}}
\toprule
Modèle & Population & Usagers & Log-loss & Brier & AP & ROC-AUC \\
\midrule
Référence constante & tous les usagers de 2024 & 117\,278 & 0,4703 & 0,1471 & 0,1793 & 0,5000 \\
Référence constante & âge renseigné seulement & 114\,859 & 0,4761 & 0,1496 & 0,1830 & 0,5000 \\
Régression logistique & tous les usagers de 2024 & 117\,278 & 0,4001 & 0,1268 & 0,4089 & 0,7699 \\
Régression logistique & âge renseigné seulement & 114\,859 & 0,4084 & 0,1294 & 0,4090 & 0,7641 \\
Arbres boostés & tous les usagers de 2024 & 117\,278 & 0,3845 & 0,1213 & 0,4629 & 0,7890 \\
Arbres boostés & âge renseigné seulement & 114\,859 & 0,3925 & 0,1238 & 0,4630 & 0,7836 \\
\bottomrule
\end{tabularx}
\sourcecalcul{notebook 08 ; mêmes prédictions dans les deux populations, sans réentraînement}
\end{table}


Retirer ces 2\,419 usagers de l'évaluation de 2024 dégrade tous les scores~: la ROC-AUC
de la régression logistique passe de 0,7699 à 0,7641, celle des arbres de 0,7890 à
0,7836. Environ \textbf{0,6 point de ROC-AUC provient donc d'une pratique
d'enregistrement}, et non d'une information sur le risque routier.

Un détail complète la lecture. L'\textit{average precision} ne bouge pratiquement pas~:
0,4089 contre 0,4090 pour la logistique, 0,4629 contre 0,4630 pour les arbres. Ces
usagers reçoivent des probabilités très basses et sont effectivement négatifs~; ils se
rangent en bas du classement sans peser sur la courbe précision-rappel. Les indicateurs
qui dépendent de toute la distribution sont donc flattés par ce groupe, alors que celui
qui se concentre sur le haut du classement ne l'est pas. Comparer deux modèles sur la
seule ROC-AUC exposerait à cette confusion.

Ce chapitre raisonne sur l'absence du seul âge, soit 2\,419 usagers en 2024. Le
chapitre~\ref{chap:test-ml} examinait l'absence de l'âge \emph{ou} du sexe, soit
2\,424 usagers~: les deux définitions se recouvrent largement sans coïncider.

\subsection{S'agit-il d'une fuite de cible~?}
Non, au sens strict~: l'information est disponible au même moment que les autres
variables et n'est pas dérivée de la gravité. Le modèle exploite une régularité
réelle et reproductible, présente aussi bien en 2023 qu'en 2024.

C'est en revanche un indicateur du \emph{processus d'enregistrement} plutôt que du
risque. Sa stabilité dépend de pratiques administratives, non d'un mécanisme routier~:
une évolution de la saisie suffirait à le faire disparaître. Les conclusions du projet
tiennent sans lui, comme le montrent les scores restreints, mais une part des
performances annoncées lui est imputable et doit être déclarée.

\section{Ce que ce chapitre permet de conclure}
Le modèle additif n'épuise pas l'information contenue dans ces onze variables. Un modèle
capable de représenter des combinaisons améliore les quatre scores, principalement sur
l'identification des cas les plus exposés. La régression logistique conserve en revanche
un avantage de lisibilité~: ses contrastes du chapitre~\ref{chap:validation-ml}
s'interprètent directement, ce que les arbres ne permettent pas sous cette forme.

Le rappel limité au seuil de 0,5 reflétait le seuil choisi, non une incapacité du modèle.
Rendre l'arbitrage explicite montre qu'un rappel de 50\,\% est atteignable au prix de
signaler un usager sur cinq, avec une règle qui se transporte d'une année à l'autre.

Enfin, une partie mesurable des performances provient d'un artefact de collecte daté et
décrit. Le signaler ne diminue pas le résultat~: cela en fixe la portée.

Aucun de ces éléments ne modifie les limites générales de l'étude. Les associations ne
sont pas causales, la population reste conditionnelle à la présence dans BAAC, les
variables ne définissent pas un instant de prédiction réel et aucune utilité
opérationnelle n'est démontrée.

\FloatBarrier
\bilan{Les arbres boostés améliorent les quatre scores à variables et protocole
identiques~; le gain porte surtout sur le haut du classement, là où le modèle additif
plafonne. L'arbitrage précision-rappel est rendu explicite et les seuils fixés sur 2023
se transportent à 2024. Environ 0,6 point de ROC-AUC provient d'une pratique
d'enregistrement apparue en 2021, désormais datée, décrite et chiffrée.}
